\def\topic{Alternative Semantics}
\ifdefined\MyHeaderIncluded
\else
\def\qwertyflag{true}

\ifdefined\MyHeaderIncluded

\else
  \newcommand{\foobar}{FooBar}
  \def\MyHeaderIncluded{true}
% if is till the end of the file

\documentclass[10pt,a4paper]{article}
\usepackage[utf8]{inputenc}
%\usepackage[german]{babel}
%\usepackage[cyrillic]{babel}
\usepackage[T1]{fontenc}
%\usepackage[T2A]{fontenc}
%\usepackage[english, russian]{babel}
\usepackage{amsmath}
\usepackage{amsfonts}
\usepackage{amssymb}
\usepackage{amsthm}
\usepackage{mathtools} %coloneqq

\usepackage{bussproofs} %

%\renewcommand{\thesubsection}{\arabic{subsection}}
%\newtheorem{thm}{Theorem}[section]
%\renewcommand{\thethm}{\arabic{subsection}.\arabic{thm}}

\usepackage{thmtools}
\usepackage{hyperref}
\hypersetup{
    colorlinks=true,
    linkcolor=blue,
    filecolor=magenta,      
    urlcolor=cyan,
    pdftitle={Overleaf Example},
    %pdfpagemode=FullScreen,
    }
%\usepackage{stix-type1}%\turnediota
\theoremstyle{definition}

\newtheorem{constructor}{Constructor}[section]
\newtheorem{iaxiom}{Axiom}
\newtheorem{axiom}{Axiom}
\newtheorem*{axiom*}{Axiom}
\newtheorem{definition}{Definition}[section]
\newtheorem*{definition*}{Definition}
\newtheorem{remark}{Remark}[section]
\newtheorem*{remark*}{Remark}
\newtheorem{recall}{Recall}[section]
\newtheorem*{recall*}{Recall}

\newtheorem{lemma}{Lemma}[section]
\newtheorem{theorem}{Theorem}[section]
\newtheorem*{theorem*}{Theorem}
\newtheorem{note}{Note}[section]
\usepackage{etoolbox}

%\newcommand*{\@Dashv}[2]{%
%  \reflectbox{$\m@th#1#2$}%
%}
\newcommand*{\separation}[3]{% 1)binder, 2)set, 3)condition
  \{#1\in #2: #3\}
}

% https://www.freecodecamp.org/news/simple-regex-tricks-for-beginners-3acb3fa257cb/
%  \\mathcal\{T\} |---> \topology
%  \\mathcal\{([^\}]*)\}   |--->  $1
\newcommand{\AND}{\ \ \&\ \ }
\newcommand{\any}{\mathrm{\ any}}
\newcommand{\borel}{\mathcal{B}}
\newcommand{\can}{\mathrm{can}}
\newcommand{\canonize}{\mathrm{canonize}}
\newcommand{\closure}{\mbox{cl}}
\newcommand{\complex}{\mathbb{C}}
\newcommand{\converges}{\longrightarrow}
\newcommand{\countable}{\mbox{\ countable}}
\newcommand{\ldefined}{\mbox{\ defined}}
\newcommand{\dom}{\mbox{dom}}
\newcommand{\Dom}{\mbox{Dom}}
\newcommand{\Done}{\noindent \textbf{Done.}\\}
\newcommand{\lexists}{\mbox{\ exists}}
\newcommand{\lforall}{\mbox{\ forall}}
\newcommand{\lunique}{\mbox{\ unique}}
\newcommand{\family}{\mathrm{\ family}}
\newcommand{\Fm}{\mathrm{Fm}}
\newcommand{\FuncSymb}{\mathrm{FuncSymb}}
\newcommand{\function}{\mathrm{\ function}}
\newcommand{\functional}{\mbox{\ functional}}
\newcommand{\IFF}{\ \Longleftrightarrow\ }
\renewcommand{\iff}{\ \longleftrightarrow\ }
\newcommand{\IMPLIES}{\ \Longrightarrow\ }
\renewcommand{\implies}{\ \longrightarrow\ }
\newcommand{\Inf}{\mathrm{Inf}}
\newcommand{\injective}{\mbox{\ injective}}
\newcommand{\injection}{\mbox{\ injection}}
\newcommand{\interior}{\mbox{int}}
\newcommand{\intersection}{\bigcap}
\newcommand{\isempty}{\mathrm{\ empty}}
\newcommand{\LB}{\mathrm{LB}}
\newcommand{\nghbset}{\mathcal{N}}
\newcommand{\nonempty}{\mathrm{\ nonempty}}
\newcommand{\Ord}{\mathrm{Ord}}
\newcommand{\power}{\mathcal{P}}
\newcommand{\pr}{\mathrm{pr}}
\newcommand{\fst}{\mathrm{fst}}
\newcommand{\snd}{\mathrm{snd}}
\newcommand{\PredSymb}{\mathrm{PredSymb}}
\newcommand{\Proof}{\noindent \textbf{Proof.}\\}
\newcommand{\ran}{\mathrm{ran}}
\newcommand{\Ran}{\mathrm{Ran}}
\newcommand{\real}{\mathbb{R}}
\newcommand{\relation}{\mathrm{\ relation}}
\newcommand{\set}{\mathrm{\ set}}
\newcommand{\singleton}{\mathrm{\ singleton}}
\newcommand{\semantics}{\mathrm{semantics}}
\newcommand{\semialgebra}{\mathcal{S}}
\newcommand{\Sup}{\mathrm{Sup}}
\newcommand{\surjective}{\mbox{\ surjective\ on\ }}
\newcommand{\Term}{\mathrm{Term}}
\newcommand{\Type}{\mathrm{Type}}
\newcommand{\topology}{\mathcal{T}}
\newcommand{\total}{\mbox{\ total\ on\ }}
\newcommand{\transitive}{\mathrm{\ transitive}}
\newcommand{\union}{\bigcup}
\newcommand{\lTrue}{\mathrm{True}}
\newcommand{\lFalse}{\mathrm{False}}
\newcommand{\lUncertain}{\mathrm{Uncertain}}
\newcommand{\lif}{\mathrm{if}}
\newcommand{\lthen}{\mathrm{then}}
\newcommand{\lelse}{\mathrm{else}}
\newcommand{\lexception}{\mathrm{exception}}
\newcommand{\lundef}{\mathrm{undef}}
\newcommand{\univocal}{\mathrm{\ univocal}}
\newcommand{\univocalat}{\mathrm{\ univocal\ at\ }}
\newcommand{\UB}{\mathrm{UB}}
\newcommand{\Un}{\mathcal{U}}
\newcommand{\Var}{\mathrm{Var}}

\newcommand{\notvarphi}{\phi}
\renewcommand{\phi}{\varphi}

\newcommand{\comment}[1]{\quad\mbox{(#1)}}
% local commands
%\preto\section{\undef\blipo} %etoolbox
\preto\subsection{\undef\blipo} %etoolbox
% https://tex.stackexchange.com/questions/105684/limit-macro-newcommand-scope-to-section
\usepackage{graphicx}
\usepackage{authblk}

\setlength\parindent{0pt} % whole doc of \noindent
\usepackage{soul}
%\usepackage{mathgl}

\fi % guard
\ifdefined\topic
\else
\def\topic{No Titel}
\fi
\title{\topic}
\author[1]{Georgy Dunaev}
\begin{document}
\maketitle
\fi
\section{[OUTDATED]Alternative semantics}

This historical proposal interprets a term by a collection of possible
classes rather than by a single class. It investigates how membership,
choice, and definedness might be expressed under that interpretation.
The label records its place in the development; the definitions below are
to be read in this local setting.

\subsection{Assignments and truth}
Let $K$ denote the conglomerate of all classes in the intended model $M$.
Both $K$ and $\power(K)$ are metatheoretic collections here; the notation
does not assert the existence of a class of all classes inside NBG or MK.
An assignment has the form
\[A:\mbox{Var}\to \power(K).\]
Thus a variable, and more generally a term, may have no candidate value,
one candidate, or several candidates. The subscripts in
$I_A^{\mathrm{Term}}$ and $I_A^{\mathrm{Fm}}$ distinguish term interpretation
from formula interpretation under the same assignment.

The proposed definedness clause is
\[
I^{\mathrm{Fm}}_A(S\downarrow)
\iff (I^{\mathrm{Term}}_A(S)\neq\emptyclass).
\]
Here definedness means nonemptiness of the candidate collection, not
uniqueness of its member. Formula interpretation is intended to remain
two-valued:
\[\mbox{Bool} \ceq \{\mbox{False}, \mbox{True}\}.\]
Satisfaction under an assignment and validity in the model are distinguished
by
\[
M,A\vDash \varphi \stackrel{meta}{\Longleftrightarrow}
I_A^{\mathrm{Fm}}(\varphi)=\mbox{True},
\]
\[
(M\vDash \varphi)
\stackrel{meta}{\Longleftrightarrow}
(\forall (A:\mbox{Assignment}),\ M,A\vDash \varphi).
\]
The quantifier in the second display ranges over the assignments admitted
by this semantics.

\subsection{Candidate membership and choice}
One proposed clause for membership, called ``weak'' in the draft, is
\[
M,A\vDash X\in Y
\quad\Longleftrightarrow\quad
\forall C\in I_A^{\mathrm{Term}}(X)\,
\forall D\in I_A^{\mathrm{Term}}(Y),\ C\in D.
\]
Here $X$ and $Y$ are arbitrary terms. Validity of $X\in Y$ additionally
quantifies over assignments. The clause compares every candidate on the
left with every candidate on the right. With ordinary universal
quantification it is vacuously true if either candidate collection is
empty; whether to add a nonemptiness condition is a further design choice.

\paragraph{A precise relational fragment.}
To isolate what this clause says, work in the metatheory with a set $D$ of
candidate objects and a binary relation $R\subseteq D\times D$. No
assumption that $D$ contains all classes is needed for this fragment. For
$S,T\subseteq D$, put
\[
R^{\forall}(S,T)\iff
\forall s\in S\,\forall t\in T,\ (s,t)\in R.
\]
The proposed membership clause is the same construction with the actual
membership relation on its candidate domain.

\begin{lemma}[Singletons, refinement, and absence]
The lifted relation has the following properties:
\begin{enumerate}
\item $R^{\forall}(\{s\},\{t\})$ holds exactly when $(s,t)\in R$.
\item If $R^{\forall}(S,T)$, $S'\subseteq S$, and $T'\subseteq T$, then
  $R^{\forall}(S',T')$.
\item $R^{\forall}(\emptyclass,T)$ and
  $R^{\forall}(S,\emptyclass)$ hold for every $S,T$.
\end{enumerate}
\end{lemma}
\Proof
In the first claim the two quantifiers range over one pair. For the second,
each pair in $S'\times T'$ is already in $S\times T$, where the relation
holds by hypothesis. In the third, there is no pair on which the relation
could fail.
\Done

This gives a conservative reading of atomic formulas on single-valued
assignments and explains why narrowing the candidate collections preserves
a true atomic membership claim. It also proves that this membership clause
alone cannot imply definedness. For example, assigning no candidates to
$X$ makes $X\in Y$ true while $X\downarrow$ is false.

If an application needs membership to imply existence of both values, it
can instead use the separately named guarded lift
\[
R^+(S,T)\iff
(S\neq\emptyclass)\land(T\neq\emptyclass)\land R^{\forall}(S,T).
\]
The singleton property remains valid, and refinement preserves truth when
the refined candidate collections remain nonempty. These conclusions follow
from the lemma and the two guards. The guarded lift is an optional variant;
the original unguarded clause is retained above.

\paragraph{Repeated occurrences and equality.}
A candidate collection records possible values but not correlations among
occurrences. If equality is lifted by the same independent-pair rule, then
\[
\operatorname{Eq}^{\forall}(S,S)\iff \forall s,t\in S,\ s=t.
\]
For nonempty $S$, this holds exactly when $S$ is a singleton: fix one
$s\in S$ and use the displayed condition to identify every other member
with it. Thus a term with two distinct candidates would fail reflexivity
under independent-pair equality. This is not an objection to the proposed
membership clause; it shows why the equality clause requires a separate
choice. Equality of candidate collections is reflexive, while evaluation
under a shared selection preserves the correlation between repeated
occurrences. Those alternatives answer different questions and must not be
silently identified.

The proposed choice principle is
\[
P(\varepsilon x,\ P)\iff \exists x,\ P(x).
\]
The letter $P$ denotes a predicate, leaving $A$ for assignments. This formula
is a hypothesis about a future descriptor interpretation. The membership
and definedness clauses alone do not establish it; the meaning of
$\varepsilon$ and of substitution into $P$ must also be specified.

\subsection{A finite model of candidate choice}
The preceding clauses admit a complete small interpretation once the
quantifier range and the descriptor are fixed. This is a local model of
a fragment, not a model of the intended class theory or its comprehension
scheme.

Fix a finite nonempty set $D$, a relation $R\subseteq D\times D$, and
a linear ordering of $\mathcal P(D)$. Use the set $\mathcal P(D)$ as
the domain of semantic objects. Variables range over candidate sets
$S\subseteq D$; a constant $c_S$ names each such set. Interpret the
binary relation symbol by $R^{\forall}$, equality by equality of candidate
sets, and $t\downarrow$ by nonemptiness of the candidate set denoted
by $t$. Give Boolean connectives and quantifiers their ordinary
classical meanings on this domain. Thus $\exists x$ ranges over
\emph{all candidate sets}, including the empty candidate set, rather
than just over individual candidates in $D$.

For a formula $\varphi$ and assignment $a$, let
\[
\mathcal W_{\varphi,a}
 =\{S\subseteq D:M,a[x\mapsto S]\vDash\varphi\}.
\]
Interpret the descriptor $\varepsilon x,\varphi$ as the first member
of $\mathcal W_{\varphi,a}$ in the fixed ordering if this set is
nonempty, and as $\emptyclass$ otherwise. The formula inside a
descriptor is a shorter expression, so these clauses define term and
formula interpretation simultaneously by recursion. No choice axiom is
needed for this finite model.

\begin{lemma}[Substitution and extensionality of finite choice]
For a term $t$, write $\llbracket t\rrbracket_a$ for its candidate-set
value. If substitution is capture-avoiding, then
\[
\llbracket r[t/z]\rrbracket_a
 =\llbracket r\rrbracket_{a[z\mapsto\llbracket t\rrbracket_a]},
\]
and the analogous equivalence holds for satisfaction of formulas.
Furthermore, formulas with the same witness collection under $a$ have
equal descriptor values under $a$.
\end{lemma}
\Proof
Use simultaneous structural induction on terms and formulas. Variables,
constants, predicates, Boolean connectives, and quantifiers obey the
usual assignment-update clauses. At a descriptor $\varepsilon x,\psi$,
rename its binder so that $x\ne z$ and $x$ is not free in $t$.
For each candidate set $S$, the induction hypothesis identifies the
truth of $\psi[t/z]$ under $a[x\mapsto S]$ with the truth of $\psi$
under $a[z\mapsto\llbracket t\rrbracket_a][x\mapsto S]$.
Here the value of $t$ is unchanged by updating $x$: a simultaneous
induction on expressions proves that values and satisfaction depend
only on free variables. Thus the two witness collections are equal.
Their first members in the fixed ordering are equal, or both collections
are empty and both defaults are empty. This proves the descriptor case.
The same reasoning proves the last assertion directly, without any
substitution: the choice rule depends on the witness collection, not
on the syntax used to describe it.
\Done

This extensionality is a property of the specified selector. A bare
axiom saying only that a descriptor chooses some witness would not force
two syntactically different descriptions with the same multiple-witness
collection to make the same choice.

\paragraph{Choice and a change of labels.}
Extensional choice need not be invariant under every relabelling of its
objects. The order used above is additional structure, and it must be
transported when the model is transported.

\begin{proposition}[No fully permutation-invariant selector]
Let $B$ be a set with at least two elements. There is no function
\[
c:\mathcal P(B)\setminus\{\emptyclass\}\longrightarrow B,
\qquad c(W)\in W,
\]
such that $c(\pi[W])=\pi(c(W))$ for every permutation $\pi$ of $B$
and every nonempty $W\subseteq B$.
\end{proposition}
\Proof
Put $b=c(B)$ and choose $b'\in B$ different from $b$. The permutation
that exchanges $b$ and $b'$ and fixes all other elements satisfies
$\pi[B]=B$. The required identity would give
$b=c(B)=c(\pi[B])=\pi(c(B))=b'$, a contradiction.
\Done

This obstruction also applies to the candidate model even if only
relabellings of the underlying $D$ are allowed. When $D$ has at least
two elements, the witness collection
$W=\{\{d\}:d\in D\}$ is invariant under every such relabelling.
A selector must choose one singleton; exchanging its member with
another member of $D$ fixes $W$ but moves that singleton. Thus even
equality of the transported witness collections cannot by itself ensure
equivariance of the selected value. This is a statement about selectors
on all nonempty collections; it does not assert that every such
collection is definable by the formulas of a given smaller language.

\begin{proposition}[Transport of ordered finite choice]
Let $(B,\leq)$ and $(B',\leq')$ be finite nonempty linearly ordered
sets, and let $h:B\to B'$ be an order-isomorphism. For each nonempty
$W\subseteq B$, put $c_{\leq}(W)=\min_{\leq}W$. Then
\[
c_{\leq'}(h[W])=h(c_{\leq}(W)).
\]
If the selectors are extended to empty collections by defaults
$b_0\in B$ and $b'_0\in B'$, the same identity holds for every
$W\subseteq B$ exactly when $b'_0=h(b_0)$.
\end{proposition}
\Proof
The element $m=\min_{\leq}W$ belongs to $W$ and satisfies $m\leq w$
for all $w\in W$. Hence $h(m)$ belongs to $h[W]$ and is below each
of its elements in the order $\leq'$. It is therefore its minimum.
For the empty collection the identity is precisely $b'_0=h(b_0)$.
\Done

For example, a bijection $f:D\to D'$ induces
$h:\mathcal P(D)\to\mathcal P(D')$, $h(S)=f[S]$.
Transport the relation by
$R'=\{(f(d),f(e)):(d,e)\in R\}$, the assignment by
$a'(x)=h(a(x))$, and each constant $c_S$ to $c_{h(S)}$.
Transport the order by declaring $h(S)\leq' h(T)$ exactly when
$S\leq T$. The default is respected because $h(\emptyclass)=\emptyclass$.
Simultaneous induction on terms and formulas now shows that term
values are carried by $h$ and formula truth is preserved. The relation
case follows by transporting every pair of candidates; quantifiers
use the bijectivity of $h$. At a descriptor, the induction hypothesis
transports its witness collection to its direct image under $h$, and
the preceding proposition transports the selected value.

Keeping the old order while permuting the objects is a different
operation: it need not preserve descriptor values. None of this
conflicts with the substitution lemma, which compares assignments
inside one fixed ordered model and identifies the witness collections
there literally. Nor does renaming a bound variable change the order
on semantic objects. Syntactic renaming, substitution, and transport
of a semantic structure are three distinct operations.

\begin{theorem}[Choice and candidate presence are distinct]
For this interpretation, with capture-avoiding substitution,
\[
M,a\vDash\varphi[(\varepsilon x,\varphi)/x]
\quad\Longleftrightarrow\quad
M,a\vDash\exists x\,\varphi.
\]
Nevertheless, the truth of either side need not imply
$(\varepsilon x,\varphi)\downarrow$.
\end{theorem}
\Proof
If there is a witness candidate set, the descriptor denotes one of the
witness sets, so substitution yields a true formula. If there is none,
$\varphi$ is false at every candidate set, including the default
$\emptyclass$, so the substituted formula is false. The substitution
step follows by induction on expressions, as in ordinary first-order
semantics; the descriptor case uses the same set of witnesses under the
two corresponding assignments.

For the last claim, take $\varphi=\neg(x\downarrow)$. Its unique
witness in $\mathcal P(D)$ is the empty candidate set. The existential
formula and the descriptor witness formula are therefore true, but the
descriptor's candidate-presence assertion is false.
\Done

\begin{corollary}[Choice with an explicit presence guard]
Define the term
\[
e_\varphi=(\varepsilon x,\ x\downarrow\land\varphi).
\]
Then
\[
e_\varphi\downarrow\iff\exists x\,(x\downarrow\land\varphi),
\qquad
 e_\varphi\downarrow\ \Longrightarrow\ \varphi[e_\varphi/x].
\]
\end{corollary}
\Proof
When a nonempty witness candidate set exists, the descriptor selects
such a set and it satisfies $\varphi$. Otherwise the witness collection
for the guarded predicate is empty and the descriptor has the empty
default, so its presence assertion is false.
\Done

In this model every term denotes a semantic object, which may itself be
an empty collection of candidates. Thus $t\downarrow$ is an internal
presence predicate, not the metatheoretic assertion that term
interpretation exists. PFOL makes a different distinction: an undefined
term has no denotation. The finite model establishes consistency of the
candidate clauses with their stated choice principle, while identifying
exactly which meaning of quantification makes that claim true. It does
not settle a semantics in which quantifiers range only over singleton
candidate sets.

\subsection{A proposed comprehension scheme}
The class-existence principle motivates a syntactic operation
$\mathrm{setify}$ on formulas. Its intended role is to restrict the relevant
variables and quantifiers to sets. A recursive definition must state how it
handles binders and free class parameters; those clauses are not yet given
in this draft.

The proposed characterization was
\[
(\varphi\in\mathrm{setify}[\mathrm{Fm}])
\stackrel{meta}{\IFF}
(\vdash (\cc{x}{\varphi}\downarrow)).
\]
The right-hand side is a claim of provability, so the equivalence requires
a specified calculus as well as semantics. In particular, the reverse
direction says that the syntactic transformation captures \emph{every}
provably defined comprehension, which is stronger than a sufficient
formation condition. The display records that intended characterization;
its two directions remain separate proof obligations.

\subsection{A precise quantifier restriction in the candidate model}
The following construction isolates one part of the preceding proposal.
It specifies a restriction of bound variables in the finite candidate
model; it does not identify that restriction with the intended
$\mathrm{setify}$ operation or interpret comprehension.

Fix the earlier finite model with semantic domain $B=\mathcal P(D)$.
Work first in its descriptor-free relational language: terms are variables
and constants $c_S$, and atomic formulas use candidate equality, the
lifted relation $R^{\forall}$, and candidate presence. Choose any subset
$E\subseteq B$ of allowed bound-variable values. Expand the language
by a fresh unary predicate $G_E$, interpreted by
\[
M,a\vDash G_E(t)\quad\Longleftrightarrow\quad
\llbracket t\rrbracket_a\in E.
\]
This is an explicitly supplied predicate interpretation, not an assertion
that $E$ is definable in the original language.

Define $M,a\vDash_E\varphi$ by keeping the existing atomic and Boolean
clauses, but letting every quantifier range over $E$. For example,
\[
M,a\vDash_E\forall x\,\psi
\quad\Longleftrightarrow\quad
\forall S\in E,\ M,a[x\mapsto S]\vDash_E\psi.
\]
Assignments and constants still have values in $B$; in particular,
free parameters are not silently restricted to $E$. This is a semantics
with restricted bound-variable ranges, not necessarily an induced
substructure with domain $E$.

Define a translation $\varphi\mapsto\varphi^E$ recursively. It fixes
atomic formulas, commutes with Boolean connectives, and has the binder
clauses
\[
\begin{split}
(\forall x\,\psi)^E&=\forall x\,(G_E(x)\implies\psi^E),\\
(\exists x\,\psi)^E&=\exists x\,(G_E(x)\land\psi^E).
\end{split}
\]
The superscript denotes a syntactic translation using the new predicate
symbol. Since that symbol has no free parameters, its inserted occurrence
introduces no new free variable and captures none from $\psi$.

\begin{proposition}[Agreement for restricted quantifiers]
\label{prop:candidate_quantifier_relativization}
For every descriptor-free formula $\varphi$ and every assignment
$a:\mathrm{Var}\to B$,
\[
M,a\vDash_E\varphi\quad\Longleftrightarrow\quad
M,a\vDash\varphi^E.
\]
This includes $E=\emptyclass$, with universal quantifiers vacuous and
existential quantifiers false in the restricted interpretation.
\end{proposition}
\Proof
Induct on the formula. Atoms have identical term values and identical
interpretations on the two sides, so the assertion holds for atoms.
Each Boolean step follows by applying its ordinary two-valued truth
operation to the equivalences supplied by induction.

For a universal quantifier, the left side means that $\psi$ holds in
the restricted interpretation under every update $a[x\mapsto S]$
with $S\in E$. By induction this is equivalent to truth of $\psi^E$
under exactly those updates. In the original domain $B$, requiring
$G_E(x)\implies\psi^E$ at every update imposes the same condition:
updates outside $E$ satisfy the implication vacuously. For an existential
quantifier, a witness update in $E$ satisfies $\psi$ in the restricted
interpretation exactly when it satisfies $\psi^E$ by induction. This
is precisely a witness in $B$ for $G_E(x)\land\psi^E$.
These arguments also handle $E=\emptyclass$, completing the induction.
\Done

\paragraph{Different restrictions give different theories.}
Take $D=\{0,1\}$ with $0\ne1$. Three possible ranges are
\[
E_{\rm all}=\mathcal P(D),\qquad
E_{\rm present}=\mathcal P(D)\setminus\{\emptyclass\},\qquad
E_{\rm single}=\{\{0\},\{1\}\}.
\]
The sentence $\forall x\,(x\downarrow)$ is false for $E_{\rm all}$
(the empty candidate set is a witness to failure) and true for
$E_{\rm present}$ and $E_{\rm single}$. The sentence
\[
\forall x\,(x=c_{\{0\}}\lor x=c_{\{1\}})
\]
is true for $E_{\rm single}$ and false for $E_{\rm present}$,
since $\{0,1\}$ is not equal to either singleton candidate set.
Thus ranging over arbitrary candidate collections, nonempty candidate
collections, and single-valued candidates are genuinely different
interpretations even when every underlying object in $D$ is a set.
None is selected by the word ``set'' alone in the original proposal.

\paragraph{Why descriptors and comprehension remain separate.}
The lemma deliberately contains no descriptor term. To extend it, one
must decide whether an $\varepsilon$ binder also restricts its witnesses
to $E$, and what happens when there is no such witness. The previous
finite choice model returns the empty candidate set for
$\varepsilon x,\bot$. This lies outside both $E_{\rm present}$ and
$E_{\rm single}$. Therefore a semantics whose term values must stay
inside either restricted domain cannot retain that default unchanged.
One can instead keep the ambient value domain $B$ while restricting
only selected quantifiers, or specify a different descriptor/default
policy; those are additional choices, not consequences of relativization.

Likewise, the chapter supplies no semantic clause for a comprehension
term $\cc{x}{\varphi}$ and no calculus for its presence assertions.
The proved translation lemma therefore does not establish either
direction of the proposed provability characterization of
$\mathrm{setify}[\mathrm{Fm}]$. It identifies a binder mechanism that
can be adopted once the intended range, parameter policy, and term
constructors are fixed, while leaving that original design question open.

\ifdefined\qwertyflag
\undef\qwertyflag
\end{document}
\fi
